\section{Evaluation: What Is Actually Measured?}
\label{sec:evaluation}
Evaluation connects representation and editability claims to observations.
The field currently reports geometric similarity, semantic image-text compatibility, aesthetic judgments, reconstruction accuracy, user-task efficiency, and editing fidelity.
These outcomes should not be combined without establishing that they measure the same construct.
Section~\ref{sec:benchmarks} compares newer benchmark contracts and their scope.

% Independent source table.
\begin{table*}[t]
\centering
\caption{Representative measurement targets and the claims they do \emph{not} directly establish. These metrics characterize different tasks, not comparable scores on one leaderboard.}
\label{tab:metrics}
\small
\begin{tabularx}{\textwidth}{@{}p{0.22\textwidth}XX@{}}
\toprule
\textbf{Measured target} & \textbf{Example protocol or metric} & \textbf{Unsupported inference without further testing}\\
\midrule
Layout fidelity/diversity & Layout FID, maximum IoU, constraint-satisfaction checks \citep{layoutdm2023}. & Correct typography, image quality, or downstream editability.\\
Text-to-vector alignment & CLIP similarity, CLIP R-Precision on rendered SVG captions \citep{vectorfusion2023}. & Path-level editing success or preserved non-target geometry.\\
Vector representation quality & Chamfer-derived reconstruction error, interpolation smoothness \citep{deepsvg2020}. & Instruction-controlled vector edits or high-quality semantic text.\\
Typography fidelity & Re-rendered reconstruction L1/PSNR \citep{derendering2021}; human legibility judgments \citep{wordasimage2023}. & Accurate arbitrary text replacement or the preservation of document structure.\\
Human workflow utility & User editing time, number of manipulations, controlled preference studies \citep{vinci2021}. & Automated edit-request success for unseen tasks.\\
Layered/raster editing & Request adherence, design/layout consistency, rendering quality \citep{mildedit2026,posteromni2026}. & Native-format layer integrity unless the artifact and validity are separately tested.\\
\bottomrule
\end{tabularx}
\end{table*}


\paragraph{Comparability conditions.}
A score is comparable across systems only when the underlying task, input conditions, output representation, test split, and measurement implementation are compatible.
For example, a layout-distribution metric cannot be compared directly with an image-distribution metric merely because both are called FID, and accurate alpha reconstruction does not measure the execution of a headline-replacement instruction.
Even within one representation, synthesis from scratch, partial completion, reconstruction, and modification of an existing artifact pose different tasks and sources of error.
We therefore interpret numerical results within their original experimental contracts and use cross-paper comparisons chiefly to establish which capabilities were actually measured.

\subsection{Generation Fidelity and Geometry}
\ac{layoutdm} evaluates discrete layout synthesis through distributional and geometry-aware measures, including \ac{fid} adapted to layout representations and maximum \ac{iou} under compatible element categories \citep{layoutdm2023}.
The conditioned generation scenarios and the quality/diversity trade-off are central to interpretation.
Such metrics concern the arrangement of design elements; even a strong score does not measure the usability of an exported creative file.

\ac{vectorfusion} evaluates prompt consistency using \ac{clip} similarity and \ac{clip} R-Precision over 128 caption prompts and analyzes the effect of vector primitive count and initialization \citep{vectorfusion2023}.
These experiments measure semantic correspondence of the rendered graphic.
By contrast, \ac{deepsvg} studies reconstruction error, interpolation smoothness, and participants' preference rankings for icon interpolations \citep{deepsvg2020}; \ac{svgformer} reports representation-learning evaluations including reconstruction, classification, and retrieval \citep{svgformer2023}.
Their benchmarks are valuable within their tasks but do not evaluate instruction-following design edits.

For layered synthesis, \ac{taue} studies 1,770 filtered MS-COCO samples using \ac{fid}, image/text alignment, and foreground/background reconstruction through \ac{psnr}, \ac{ssim}, and \ac{lpips} (Section 4.2 and Table 2).
Its Figure 6 shows layout and background-replacement demonstrations, but Section 6 acknowledges pixel-level foreground-preservation limitations, and its reconstruction scores do not measure held-out edit-instruction success.
\ac{tkgdm} evaluates chroma-key images generated for 3,000 MAGICK-derived prompts with \ac{fid}, mask-based FID, and image/text \ac{clip} alignment scores (Section 5.2 and Table 1), rather than measuring operations on native editing objects.
Crucially, the reported SDXL FID is 41.81 for \ac{tkgdm} versus 29.34 for LayerDiffuse in the same table, so the paper does not establish universal superiority over fine-tuned layer-generation approaches.
Its 100-participant two-alternative preference study (Section 5.5) measures perceived foreground quality and text alignment, not controlled graphic-design edit success.
In \ac{alfie}, the fraction of 3,000 generated illustrations with an empty four-pixel border on all sides rises from 53.10\% under a prompt-suffix baseline to 96.50\% with subject centering (Table 1); this is a border-based proxy for uncropped assets, not a measurement of semantic editability.
A separate ten-person study on 100 images prefers Alfie's alpha mattes over ViTMatte approximately 63\% of the time (Section 5.1 and Figure 6), while Figure 7 illustrates downstream collage compatibility without quantitative instruction-edit evaluation.
The authors explicitly omit \ac{fid} and \ac{kid} for \ac{rgba} outputs because an appropriately matched transparent-image reference set is unavailable (Section 5), so the missing scores must not be inferred from CLIP-S or border statistics.
In contrast, \ac{layeringdiff} measures 572 generated foreground/background/composite prompt triplets with \ac{fid}, \ac{kid}, \ac{clip} score and a detector-derived \ac{fgmiou} metric (Section 4.1 and Table 1).
Its composite \ac{fid} of 121.05 improves on LayerDiffuse T2L's 134.51 under the reported setup, while its foreground \ac{fid} of 133.76 is worse than the baseline's 127.14, so the work does not dominate every layer-wise metric.
Table 3 measures image quality and prompt alignment from 24 participant ratings, and Table 4 reports reconstruction error on 1,000 synthetic composites, not instruction-based post-export edit performance.
For background recovery, Table 4 reports an improved \ac{lpips} of 21.80 versus 28.73 for the Germer baseline, but worse MSE of 17.52 versus 12.76, demonstrating that perceptual and pixel-reconstruction metrics answer distinct questions.
Its Section 5 also reports alpha-estimation errors and shadows remaining in background layers as failure modes that can constrain later compositing.
\ac{uniworlddesign} evaluates I2L on 512 held-out Crello designs with four predicted layers, aligned to source stacks by LayerD before scoring (Section 4.1 and Table 1).
Against Qwen-Image-Layered, per-layer RGB L1 decreases from 0.2014 to 0.1264 and Alpha Soft IoU increases from 0.5454 to 0.7325, but these are reconstruction and alpha-overlap measures rather than executed edit requests.
The Stable-Layers-derived structural proxy \emph{Bad Layers} declines from 1.69 to 1.32 per design (Table 2), while five-axis VLM judging raises the total from 17.60 to 20.43 out of 25 but lowers alpha cleanliness from 3.33 to 2.90 (Table 3 and Appendix B).
The VLM sees the input, recomposition and predicted-layer contact sheet, not a human-operated editor performing a controlled downstream text or object change.
For T2RGBA, 512 design-element prompts yield a mean CLIP Score of 33.03 versus OmniAlpha's 31.00, whereas OmniAlpha has better FID (87.86 versus 117.14) and white-composite LPIPS (0.462 versus 0.537) under their respective setup (Table 4).
The paper reports alpha-edge and dense-typography failures (Section 5), limiting any inference from layer separability to unrestricted native document editing.

\subsection{Appearance, Text, and Reconstruction}
Typography introduces additional evaluation dimensions: correct character identity, readability, style consistency, and editable text parameters.
\ac{derendering} reports reconstruction L1/\ac{psnr} and illustrates text or styling changes using recovered renderer parameters \citep{derendering2021}.
Those numerical measurements indicate reconstruction fidelity, not a rate of successful edit instructions.
\ac{wordasimage} includes a perceptual study measuring whether semantic concepts remain recognizable, letters remain legible, and typography resembles the input font \citep{wordasimage2023}.
The same dimensions should not be conflated with free-form replacement of textual content.
\ac{postertextbench} evaluates targeted raster-image edits, not only text rendering, with OCR-based sentence accuracy and normalized edit distance for insertion and modification, masked \ac{lpips} for non-target preservation, and VLM instruction/visual-quality scoring (Section 4 and Table 1).
In patch addition, \ac{postertext} obtains sentence accuracy 0.763, normalized edit distance 0.899 and masked \ac{lpips} 0.039, versus Nano Banana 2's 0.500, 0.635 and 0.064.
For patch modification, its sentence accuracy 0.771 and normalized edit distance 0.902 are lower than Nano Banana 2's 0.819 and 0.907, even as masked \ac{lpips} improves to 0.028 versus 0.034.
For deletion, PosterText reports masked \ac{lpips} 0.040 and VLM instruction-following 4.70, while Nano Banana Pro has higher VLM visual quality 4.91 versus 4.78.
Table 2's full-training sentence accuracy 0.767 and masked \ac{lpips} 0.036 aggregate different subsets of editing tasks rather than describing a single operation.
These image-level results demonstrate measured instruction-guided raster editing, not human editing time, export validity or preservation of native font and vector objects.

\paragraph{Color, geometry, and typography are different revision targets.}
For icon colorization, \citet{li2022iconcolorization} synthesize flat-color pixels from sketches and style conditions, which assesses image colorization rather than vector-layer edit validity.
By contrast, \citet{chen2023primitiveassembly} predict editable parametric primitives, and \citet{gehlaut2026vectorark} reconstruct rounded-polygon vector graphics from raster inputs; both deliver structured geometry, but neither task by itself scores semantic text replacement.
For typographic imagery, \citet{wang2022textlogo} predict glyph arrangements from linguistic and visual cues, whereas \citet{wang2023anythingtoglyph} synthesize artistic letter appearances.
Aesthetics++ \citep{kong2023aestheticspp} instead ranks design-principle-guided revisions of existing slide designs according to learned human preferences.
These examples should not share a single editability or quality score: pixel colorization, primitive manipulation, glyph design, and preference-ranked slide refinement have different inputs, outputs and judged operations.

\paragraph{Visual-text placement and constraint-aware design.}
Content-conditioned poster generation ranges from adversarial background-aware layout \citep{zheng2019contentgan} through element-aware positioning \citep{li2024designelementposter}, density-based layout conditioning \citep{hsu2023densitylayout}, and priority-sensitive control \citep{yang2026priorityposter}.
Text placement over natural images \citep{li2022smarttext} and staged poster layout inference \citep{chai2023twostageposter} further change the target objects and aesthetic objectives.
The relevant interaction contracts also include mobile-shopping scenarios \citep{jing2023mobilshopping}, multi-constraint interactive layout \citep{chen2024iris}, and user-guided image-aware placement \citep{xu2025imageaware}.
Advertising creative generation \citep{wang2022creagan} and retrieval-augmented agentic composition \citep{forouzandehmehr2025calrag} optimize different production pipelines; the ability to request another layout is not equivalent to measured preservation of unrelated objects after editing a saved design.

\paragraph{Verification gaps despite editable design programs.}
CreatiPoster \citep{zhang2025creatiposter} produces an editable JSON program containing text and asset layers and supports canvas-conditioned completion and re-layout (Sec.~3.4 and Table~1 of the source).
Its 90-case test collection separates 45 prompt-only cases, 39 single-asset cases, and six multi-asset cases, with VLM rubric scores for layout, color, style, and compliance and ten volunteers' overall ratings.
The asset-conditioned comparison excludes the six multi-asset cases, leaving 84 cases; its authors report a higher human overall mean for CreatiPoster-S than Canva (2.82 versus 2.33), but the Canva layout score is higher in that setting (2.84 versus 2.41 for CreatiPoster-S).
Thus evidence for an editable source protocol and preference-based image quality should be presented separately from measured post-export edit instruction success.

\subsection{Editability Affordance Versus Editing Performance}
Three complementary examples clarify the evidence distinction.
Vinci uses a controlled within-subject study for poster resizing and reports the time and number of editing operations with and without feedback support \citep{vinci2021}.
That is direct evidence about \emph{human workflow efficiency}, not about autonomous editing.
\ac{editablevisual} shows that output \ac{html}/\ac{css} can expose separate \ac{dom} objects and editable text in a set of case studies \citep{editablevisual2026}.
This demonstrates a useful workflow and format affordance but does not by itself constitute a randomized user study or an instruction-edit accuracy estimate.
\ac{iconshop} demonstrates \ac{svg} icon-path manipulation while reporting \ac{fid}, \ac{clip}-based scores, and icon quality judgments for its synthesis task \citep{iconshop2023}.
The quantitative evaluation primarily concerns synthesis, not editing.

Dedicated editing benchmarks merit distinct treatment.
\ac{svgeditbench} directly evaluates \ac{llm}-driven \ac{svg} code edits on curated instructions \citep{svgeditbench2024}, while \ac{vectorgym} extends evaluation to four tasks, including expert-authored complex \ac{svg} edits, sketch conversion, text-based vector generation, and captioning \citep{vectorgym2026}.
These are more appropriate counterexamples to any blanket claim that quantitative vector-editing evaluation is missing.
\ac{mildedit} describes evaluation axes for instruction following, design layout, aesthetics, and text on multi-layer documents \citep{mildedit2026}.
\ac{posteromni} combines local editing and global poster construction, so evidence must be separated by task and output representation \citep{posteromni2026}.
A paper can evaluate image-editing ability even when it delivers flattened raster output.
The appropriate criticism is not a lack of object-level export per se, but whether the representation and evaluation support the proposed downstream claim.

\paragraph{Generation feedback versus editing performance.}
\ac{introsvg} refines SVG code through rendered-image critique and \ac{diagramagent} generates revisable diagrams, but neither output affordance by itself establishes post-export protected-object preservation.

\subsection{An Operational Evaluation Protocol}
Figure~\ref{fig:protocol} shows the separation of target modification, protected content, artifact validity, visual quality, and effort.
We recommend reporting editing performance as a vector of outcomes, rather than a single unqualified quality score.
For an input artifact $D$, edit instruction $u$, and target set $T$, a system returns $D'$.
At minimum, evaluation should address:
\begin{equation}
\mathbf{m}(D,u,D') =
\left[
m_{\mathrm{target}},\
m_{\mathrm{preserve}},\
m_{\mathrm{valid}},\
m_{\mathrm{visual}},\
m_{\mathrm{effort}}
\right].
\label{eq:metrics}
\end{equation}
Here target adherence tests the requested change; preservation checks non-target elements, structure, and content; validity tests whether the native artifact still opens and renders; visual quality assesses the edited output; and effort quantifies human or computational cost.
These are \emph{recommended evaluation dimensions}, not a claim that a common benchmark already provides all five.

A protocol should publish the initial design and target edit, the required output format, permitted tools, allowed untouched regions, failure handling, and the evaluation judge.
It should also distinguish an edit that preserves visual appearance by flattening the design from one that preserves editability.
For automated judges, human calibration, judge variance, and adversarial failure cases are necessary, especially when aesthetics are being rated.


\paragraph{Structural layout distances and color-assistance criteria.}
\ac{ltsim} measures layout similarity through transportation of labeled and geometric elements, and therefore asks a different question from a learned-feature distributional distance such as layout \ac{fid}.
InfoColorizer studies interactive palette recommendation and designer-facing controllability for information graphics \citep{infocolorizer2021}.
\citet{qiu2023multipalette} instead predict masked colors across image, SVG, and text palettes in vector graphic documents.
On 2,278 held-out Crello designs, their multi-palette model reports top-1 masked-color accuracy of 0.36 against 0.23 without segment embeddings and 0.03 for a Word2Vec-based comparison (Table 2).
Their qualitative evaluation involves 84 participants, and their interactive recoloring system was separately discussed with 12 professional designers (Sections 4.2--4.3).
These results demonstrate color prediction and designer-facing recoloring, but do not establish the success rate of arbitrary post-export semantic editing requests.
The top-1 accuracy falls from 0.36 for one masked color to 0.17 for five (Table 5), indicating a limitation of simultaneous multi-color recommendations.
Color agreement, geometric matching, and task completion should be reported as independent evaluation dimensions, particularly when a design system supports human adjustments.

\subsection{Human Preferences and Design-Principle Judges}
Graphic-design evaluation must distinguish aesthetic preference, layout-rule compliance, and whether a model can identify a design fault.
\citet{haraguchi2024gptdesign} compare GPT-based judgment of design principles against heuristic scores and human assessments rather than treating a \ac{vlm}'s score as ground truth.
\ac{designometer} instead learns to evaluate and suggest refinements to graphic-design compositions, whereas \ac{designsense} collects human preferences for layout transformations and trains an evaluation signal in that domain.
\ac{posterreward} focuses specifically on the typography, layout, and visual quality of generated posters.
These studies provide complementary evidence for evaluation-model behavior and do not demonstrate the accuracy of a fully autonomous editing agent.
